\chapter{Формы логической детерминации власти}

\section{Тезис: власть в себе — потенция (Potentia)}

На первом уровне власть существует как «власть в себе» — как чистая потенциальность, ещё не реализованная, но уже содержащая возможность господства. Это скрытая, незаметная власть — её ещё не видно, но она уже присутствует как готовность действовать, как напряжённая возможность.

Такова власть до всякого действия, до формального закрепления, до признания. Она — как энергия, сжатая в ядре, ещё не проявленная, но уже способная определять. Её сила в её молчании.

Именно в этом виде существует харизма до лидерства, сила до командования, влияние до приказа. Это уровень, на котором власть не требует института, закона или статуса — она просто *есть*, как весомость, как смысловая гравитация, как бытие возможности.

В этой форме власть абсолютна, но неопределённа. Она ещё не выделена, не опосредована, не измерена — но она уже внутренняя необходимость: субъект знает, что может повелевать, но ещё не делает этого.

\section{Антитезис: власть как определённость — Potestas}

Власть, будучи потенциальной, не существует во внешнем мире до тех пор, пока не приобретёт форму. Эта форма — и есть *Potestas* — власть как определённое, институциональное, различимое. Здесь власть становится объектом: она закреплена, оформлена, видима.

*Potestas* — это власть в её оформленной реальности: в виде закона, титула, устава, приказа. Она действует не как возможность, а как структура. Здесь уже не просто «могу», а «имею право». Сила обретает рамку, канал, масштаб, и потому становится надличной.

Но вместе с определённостью приходит ограничение. Институционализированная власть теряет свою внутреннюю напряжённость, свою чистую волю. Она становится воспроизводимой, передаваемой, делегируемой. Теперь ею может обладать не тот, кто способен, а тот, кто поставлен.

В этом виде власть переходит от духа к форме. Она теряет абсолютность, но обретает действие. Она больше не скрыта — но уже и не свободна: её обрамляют законы, процедуры, традиции. Она становится вещью, функцией, должностью.

Так *Potestas* противостоит *Potentia*: одна свободна, но неоформлена; вторая оформлена, но подчинена структуре.

\section{Синтез: власть как становление формы — Autoritas}

Истинная власть возникает не как чистая возможность (*Potentia*) и не как институциональная структура (*Potestas*), а как становление — как живое соединение силы и формы. Это — *Autoritas*: власть, обладающая формой, но не мёртвой; обладающая силой, но не хаотичной.

*Autoritas* — это власть, в которой потенция актуализируется через форму, а форма остаётся носителем духа. Она уже признана, но всё ещё дышит. Это не просто титул, а легитимность; не просто команда, а смысл. *Autoritas* невозможна без основания в способности, но и не существует вне признанной структуры.

В этом синтезе проявляется власть как органическая: она не навязана извне, но принята как естественная. Она не ломает, а оформляет. Она обладает способностью действовать и быть признанной в этом действии как необходимая. В этом заключается её устойчивость.

*Autoritas* — это не просто быть в силе, и не просто быть в праве, а быть тем, кто соединяет оба начала в форме, которая есть нечто большее, чем сумма. Это и есть власть, способная продлевать себя во времени, преодолевать кризисы и восстанавливать легитимность.

Таким образом, власть как становление формы — не застывшая структура и не стихийная энергия, а живая логика духа, обретшая устойчивость.
